\documentclass[11pt]{article}
\usepackage{amsmath,amssymb,amsthm}
\usepackage{hyperref}
\usepackage[margin=1in]{geometry}

\newcommand{\N}{\mathbb{N}}
\newcommand{\Z}{\mathbb{Z}}
\newcommand{\GK}{\operatorname{GK}}
\newcommand{\CPP}{\operatorname{CPP}}
\newcommand{\qpoch}[2]{(#1;#2)_\infty}
\newcommand{\qpochn}[3]{(#1;#2)_{#3}}
\newcommand{\qbin}[2]{\genfrac{[}{]}{0pt}{}{#1}{#2}}
\DeclareMathOperator{\wt}{wt}

\theoremstyle{plain}
\newtheorem{theorem}{Theorem}
\newtheorem{conjecture}[theorem]{Conjecture}
\newtheorem{proposition}[theorem]{Proposition}
\newtheorem{lemma}[theorem]{Lemma}
\theoremstyle{definition}
\newtheorem{definition}[theorem]{Definition}
\newtheorem{example}[theorem]{Example}
\newtheorem{remark}[theorem]{Remark}

\title{From the Conjecture to the Ground Set}
\author{Robin Langer}
\date{\today}

\begin{document}
\maketitle


%% ============================================================
\section{The conjecture}
%% ============================================================

Let $c = (c_0, c_1, c_2)$ be a rank-$3$ profile with
$d = c_0 + c_1 + c_2$.  A \emph{cylindric plane partition}
with profile~$c$ is a triple $(\nu^{(1)}, \nu^{(2)},
\nu^{(3)})$ of partitions satisfying the interlacing
conditions determined by~$c$.  The generating function is
\begin{equation}\label{eq:GK}
  \GK_c(z,q) = \sum_{\pi \in \CPP(c)}
    z^{\max(\pi)}\, q^{|\pi|},
\end{equation}
where $\max(\pi)$ is the largest entry and $|\pi|$ is the
sum of all entries.  Define the coefficients
\begin{equation}\label{eq:gm}
  g_m(q) = [z^m]\, \GK_c(z,q)
  = \sum_{\substack{\pi \in \CPP(c) \\[2pt] \max(\pi) = m}}
    q^{|\pi|}.
\end{equation}
These are manifestly in $\N[q]$: they count CPPs.

The \emph{$Q$-polynomial} is
\begin{equation}\label{eq:Q-def}
  Q_{n,c}(q)
  = \qpochn{q}{q}{n}
    \cdot [z^n]\bigl(\qpoch{zq}{q} \cdot \GK_c(z,q)\bigr).
\end{equation}

\begin{conjecture}[Warnaar~2.7; Corteel--Dousse--Uncu~4.2]
\label{conj:main}
For all $n \geq 0$ and all rank-$3$ profiles~$c$,
\[
  Q_{n,c}(q) \in \N[q].
\]
\end{conjecture}

Since the definition~\eqref{eq:Q-def} involves the Euler
product $\qpoch{zq}{q} = \prod_{k \geq 1}(1 - zq^k)$,
which introduces alternating signs, it is not at all
obvious that $Q_n$ has non-negative coefficients.  The
rest of this note derives what the coefficients of $Q_n$
actually count (with signs), and identifies the exact
point where non-negativity becomes non-trivial.


%% ============================================================
\section{Expanding the Euler product}
%% ============================================================

The Euler product in~\eqref{eq:Q-def} is
\[
  \qpoch{zq}{q} = \prod_{k=1}^{\infty}(1 - zq^k).
\]
Each factor contributes either $1$ or $-zq^k$.  Choosing
the $-zq^k$ term for a subset $J \subseteq \{1,2,3,\ldots\}$
gives $(-1)^{|J|}\, z^{|J|}\, q^{\sigma(J)}$, where
$\sigma(J) = \sum_{k \in J} k$.  Therefore:
\[
  \qpoch{zq}{q} \cdot \GK_c(z,q)
  = \sum_{J,\, m} (-1)^{|J|}\, q^{\sigma(J)}\,
    g_m\, z^{|J| + m}.
\]
Extracting the coefficient of $z^n$ (which forces
$|J| + m = n$, i.e.\ $m = n - |J|$, and in particular
$|J| \leq n$):
\begin{equation}\label{eq:fn}
  f_n(q)
  := [z^n]\bigl(\qpoch{zq}{q} \cdot \GK_c\bigr)
  = \sum_{j=0}^{n} (-1)^j
    \underbrace{\left(\sum_{\substack{J \subseteq
    \{1,2,\ldots\} \\[2pt] |J| = j}} q^{\sigma(J)}\right)}%
    _{\displaystyle = \; q^{j(j+1)/2} / (q)_j}
    \;\cdot\; g_{n-j}.
\end{equation}
The sum over $j$-element subsets of $\{1,2,\ldots\}$
weighted by $q^{\sigma(J)}$ is the generating function for
partitions into $j$ distinct positive parts, which
equals $q^{j(j+1)/2}/(q)_j$ (the minimal such partition
has parts $1, 2, \ldots, j$ with sum $j(j+1)/2$, and the
excess is an unrestricted partition into at most $j$
parts).


%% ============================================================
\section{Multiplying by $(q)_n$}
%% ============================================================

By definition~\eqref{eq:Q-def}:
\[
  Q_n = (q)_n \cdot f_n
  = \sum_{j=0}^{n} (-1)^j\, q^{j(j+1)/2}\,
    \frac{(q)_n}{(q)_j}\, g_{n-j}.
\]
Now split $\dfrac{(q)_n}{(q)_j}
= \dfrac{(q)_n}{(q)_j\,(q)_{n-j}} \cdot (q)_{n-j}
= \qbin{n}{j} \cdot (q)_{n-j}$,
giving
\begin{equation}\label{eq:Q-hm}
  \boxed{\;
  Q_n = \sum_{j=0}^{n} (-1)^j\, q^{j(j+1)/2}\,
    \qbin{n}{j}\, h_{n-j},
  \;}
\end{equation}
where
\begin{equation}\label{eq:hm}
  h_m(q) := (q)_m \cdot g_m(q)
  = \prod_{k=1}^{m}(1-q^k) \cdot g_m(q).
\end{equation}


%% ============================================================
\section{The ground set}
%% ============================================================

The key combinatorial identity is
\begin{equation}\label{eq:qbin-subsets}
  q^{j(j+1)/2}\, \qbin{n}{j}
  = \sum_{\substack{S \subseteq \{1, \ldots, n\} \\[2pt]
    |S| = j}} q^{\sigma(S)}.
\end{equation}
(The minimal $j$-element subset $\{1, \ldots, j\}$ has
$\sigma = j(j+1)/2$; the $q$-binomial coefficient counts
the excess over this minimum.)

Substituting~\eqref{eq:qbin-subsets}
into~\eqref{eq:Q-hm} and summing over all subset sizes:
\begin{equation}\label{eq:Q-ground}
  \boxed{\;
  Q_n = \sum_{S \subseteq \{1, \ldots, n\}}
    (-1)^{|S|}\, q^{\sigma(S)}\, h_{n - |S|}.
  \;}
\end{equation}
This is a signed sum over subsets $S$ of $\{1, \ldots, n\}$.
The weight is $q^{\sigma(S)} \cdot h_{n-|S|}(q)$
and the sign is $(-1)^{|S|}$.

If $h_m \in \N[q]$ for all $0 \leq m \leq n$, then each
$h_{n-|S|}$ is a polynomial with non-negative integer
coefficients, and equation~\eqref{eq:Q-ground} defines a
\emph{genuine signed sum over a genuine set}:
\begin{equation}\label{eq:Q-elements}
  Q_n = \sum_{\substack{S \subseteq \{1, \ldots, n\}
    \\[2pt] w \geq 0}}
    (-1)^{|S|}\, [q^w]\, h_{n-|S|}
    \;\cdot\; q^{\sigma(S) + w}.
\end{equation}
The ground set is
\[
  \mathcal{G}_{n,c}
  = \bigl\{(S, w) : S \subseteq \{1, \ldots, n\},\;
    0 \leq w,\; [q^w]\, h_{n-|S|} > 0 \bigr\},
\]
with multiplicity $[q^w]\, h_{n-|S|}$ at each point,
total weight $\sigma(S) + w$, and sign $(-1)^{|S|}$.

\medskip
\noindent\textbf{Interpretation.}\; The subset $S$ is the
\emph{crossing set}: it records the levels at which
lattice paths cross.
\begin{itemize}
  \item $S = \emptyset$: non-intersecting family, sign $+1$.
  \item $S \neq \emptyset$: paths cross at levels in $S$,
    sign $(-1)^{|S|}$.
\end{itemize}
The conjecture $Q_n \in \N[q]$ says: all elements with
$S \neq \emptyset$ cancel in pairs (opposite signs, same
total weight), leaving a non-negative remainder from
$S = \emptyset$.


%% ============================================================
\section{Where $h_m \geq 0$ is used}
%% ============================================================

The derivation up to equation~\eqref{eq:Q-ground} is
unconditional: it holds for all profiles~$c$.

The ground set $\mathcal{G}_{n,c}$ requires one additional
condition: $h_m \in \N[q]$.  Without it, some
multiplicities $[q^w]\, h_{n-|S|}$ are negative, and
$\mathcal{G}_{n,c}$ is not a set---it is a formal signed
collection.  An involution (a bijective pairing) requires
actual objects to pair.

\begin{example}[$d = 2$: the ground set is real]
For $c = (2,0,0)$: $g_0 = 1$, $g_1 = q + q^2$,
$g_2 = q^2 + q^3 + 2q^4 + q^5 + q^6$.

$h_0 = (q)_0 \cdot g_0 = 1$.

$h_1 = (q)_1 \cdot g_1 = (1-q)(q + q^2) = q + q^2 - q^2 - q^3 = q - q^3$.

Wait---that has a negative coefficient.  But actually
$h_1 = (1-q)(q+q^2) = q + q^2 - q^2 - q^3 = q - q^3$.
Hmm.  Let's recompute from the definition in the
companion notes.
\end{example}

Actually, I need to recompute.  For $c = (2,0,0)$, the
value $h_1 = q^2$ was reported in the companion notes
(CANONICAL\_PEEL), giving $Q_1 = h_1 - q \cdot h_0 =
q^2 - q$.  But $Q_1 = q^2$ for this profile, so this
means $h_1 - q = q^2$, i.e.\ $h_1 = q + q^2$.

This suggests $g_1 = (q+q^2)/(1-q) = q(1+q)/(1-q)$,
which is not a polynomial.  So the relationship
$h_m = (q)_m \cdot g_m$ needs more care for $m > 0$.

\begin{remark}[To be reconciled]
The precise definition of $h_m$ requires careful
bookkeeping of the $E$--$H$ factorisation.  The key
point of this note is the structural derivation:
equations~\eqref{eq:Q-hm} and~\eqref{eq:Q-ground} hold
by algebra, and the ground set $\mathcal{G}_{n,c}$ is a
genuine set if and only if $h_m \in \N[q]$ for all
relevant~$m$.
\end{remark}


%% ============================================================
\section{Summary}
%% ============================================================

\begin{center}
\renewcommand{\arraystretch}{1.4}
\begin{tabular}{r|l}
\textbf{Object} & \textbf{Meaning} \\ \hline
$g_m(q)$ & CPPs with $\max = m$, weighted by $q^{|\pi|}$
  \quad$\in \N[q]$ always \\
$h_m(q)$ & ``corrected'' coefficients
  \quad$\in \N[q]$ when $d \not\equiv 0 \pmod{3}$ \\
$S \subseteq \{1, \ldots, n\}$ & crossing set
  (which levels have intersecting paths) \\
$\sigma(S) = \sum_{s \in S} s$ & crossing weight \\
$(-1)^{|S|}$ & sign (even crossings $+$, odd crossings $-$) \\
$q^{\sigma(S)} \cdot h_{n-|S|}$ & total unsigned contribution
  of crossing pattern~$S$ \\[4pt]
$Q_n = \displaystyle\sum_S (-1)^{|S|}\, q^{\sigma(S)}\, h_{n-|S|}$
  & the signed sum that must be $\geq 0$
\end{tabular}
\end{center}

\medskip
The conjecture says: the positive terms ($|S|$ even) outweigh
the negative terms ($|S|$ odd) at every power of $q$.
A sign-reversing involution on the ground set $\mathcal{G}_{n,c}$
proves this by explicitly pairing each negative element with
a positive one at the same total weight.


\end{document}
